% TikZ 图：家庭电路与安全保护
% Day 24：交变电流、变压器与电磁应用
\begin{tikzpicture}[scale=0.88,>=Latex,line cap=round,line join=round,
  every node/.style={font=\small}]
  \tikzset{
    wire/.style={thick},
    live/.style={thick,red!80!black},
    neutral/.style={thick,blue!70!black},
    earth/.style={thick,green!45!black},
    device/.style={draw,rounded corners=3pt,fill=gray!8,minimum width=1.35cm,minimum height=0.76cm},
    protector/.style={draw,rounded corners=3pt,fill=orange!10,minimum width=1.55cm,minimum height=0.82cm}
  }

  \node[draw,rounded corners=3pt,fill=blue!6,minimum width=1.35cm,minimum height=1.15cm] (grid) at (0,0.70) {用户端};
  \node[protector] (breaker) at (2.25,1.30) {空气开关};
  \node[protector] (rcd) at (4.55,0.70) {漏电保护};
  \node[device] (load) at (7.00,0.70) {用电器};

  \draw[live,->] (grid.east |- 0,1.30) -- (breaker.west);
  \draw[live,->] (breaker.east) -- (rcd.west |- 0,1.30);
  \draw[live,->] (rcd.east |- 0,1.30) -- (load.west |- 0,1.30);

  \draw[neutral,->] (load.west |- 0,0.10) -- (rcd.east |- 0,0.10);
  \draw[neutral,->] (rcd.west |- 0,0.10) -- (grid.east |- 0,0.10);

  \node[red!80!black,above] at (1.10,1.30) {火线 $L$};
  \node[blue!70!black,below] at (1.25,0.10) {零线 $N$};

  \draw[thick,fill=white] (6.35,-1.15) rectangle (7.65,-0.42);
  \node at (7.00,-0.80) {金属外壳};
  \draw[wire] (load.south) -- (7.00,-0.42);
  \draw[earth,->] (7.00,-1.15) -- (7.00,-2.05);
  \draw[earth] (6.55,-2.05) -- (7.45,-2.05);
  \draw[earth] (6.70,-2.22) -- (7.30,-2.22);
  \draw[earth] (6.86,-2.38) -- (7.14,-2.38);
  \node[green!45!black,right] at (7.13,-1.70) {保护地线 $PE$};

  \draw[red!80!black,dashed,->] (6.55,0.28) -- (6.72,-0.42);
  \draw[green!45!black,dashed,->] (6.72,-0.42) -- (7.00,-1.15);

  \node[font=\scriptsize,align=center] at (2.25,2.15) {总电流过大\\自动断开};
  \draw[orange!80!black,->] (2.25,1.78) -- (2.25,1.54);

  \node[font=\scriptsize,align=center] at (4.55,-1.08)
    {正常：$I_L=I_N$\\漏电：$I_L\ne I_N$};
  \draw[orange!80!black,->] (4.55,-0.42) -- (4.55,0.22);

  \node[font=\scriptsize,align=center] at (7.00,1.95)
    {电能转化为光、热\\或机械能};
\end{tikzpicture}
